\problema{Binary Tree Maximum Path Sum}{124}{src/02-arboles/124-max-path-sum.cpp}{H}

Nos dan un árbol binario cuyos valores pueden ser negativos. Un \emph{camino}
es cualquier secuencia de nodos donde cada par consecutivo está conectado por
una arista, sin repetir ningún nodo. El camino \textbf{no} necesita pasar por
la raíz ni terminar en una hoja: puede empezar y terminar en cualquier nodo.
Queremos la suma máxima de valores que puede lograr algún camino.

\paragraph{Intuición.} Párate en cada nodo y pregúntate dos cosas distintas.
Primero: "si quiero seguir subiendo este camino hacia mi padre, ¿cuánto es lo
más que puedo aportar bajando por una sola de mis dos ramas?" (subir por las
dos ramas a la vez no tiene sentido: el camino se bifurcaría). Segundo: "si mi
camino ya no va a subir más y va a \emph{doblar} justo aquí, usando mis dos
ramas a la vez, ¿cuánto sumaría?". La primera pregunta es lo que le reportas a
tu padre; la segunda solo sirve para actualizar, en silencio, la mejor
respuesta que hemos visto en todo el árbol.

\begin{center}
\ttfamily
\begin{tabular}{ccc}
 & -10 & \\
9 & & 20\\
 & 15 & 7\\
\end{tabular}
\end{center}

\noindent En este árbol el camino óptimo es $15 \to 20 \to 7 = 42$: conviene
ignorar por completo la raíz $-10$ y la rama de $9$, porque sumarlas solo
empeoraría el resultado.

\paragraph{Solución.} Hacemos DFS recursivo tipo post-order con una función
auxiliar \texttt{gananciaMaxima(nodo)} que devuelve, para el subárbol de ese
nodo, la ganancia máxima que aporta a un camino que sigue subiendo hacia el
padre (una sola rama):

\begin{center}
\ttfamily
gananciaIzq = max(0, gananciaMaxima(nodo->left))\\
gananciaDer = max(0, gananciaMaxima(nodo->right))\\
devuelve = nodo->val + max(gananciaIzq, gananciaDer)
\end{center}

\noindent El \texttt{max(0, ...)} descarta una rama si su ganancia es
negativa: es mejor no extender el camino por ahí, porque sumar un número
negativo solo lo empeora.

Es crucial no confundir lo que la recursión \emph{devuelve} (una sola rama,
porque un camino que sube hacia el padre solo puede continuar por un lado)
con lo que actualiza la \emph{respuesta global} (las dos ramas a la vez, para
el camino que dobla justo en ese nodo y ya no puede crecer más hacia arriba):

\begin{center}
\ttfamily
candidato = nodo->val + gananciaIzq + gananciaDer\\
mejorGlobal = max(mejorGlobal, candidato)
\end{center}

\noindent Recorremos en post-order (primero los hijos, luego el nodo actual)
para que cada nodo ya conozca la ganancia de sus dos subárboles antes de
calcular su propio candidato. El código completo (abre el archivo con clic en
el título):

\lstinputlisting{src/02-arboles/124-max-path-sum.cpp}

\complejidad{$O(n)$}{$O(h)$}

\paragraph{Notas de entrevista (Amazon).} El error más común es devolver
\texttt{val + gananciaIzq + gananciaDer} hacia el padre en vez de
\texttt{val + max(gananciaIzq, gananciaDer)}: eso rompería la invariante de
que un camino no se bifurca. Otro punto clave: la respuesta debe poder ser
negativa (por ejemplo un árbol de un solo nodo con valor negativo), así que
\texttt{mejorGlobal} debe inicializarse en \texttt{INT\_MIN} y nunca forzarse
a 0. Casos borde: nodo único, todos los valores negativos, y una rama
negativa que conviene podar con \texttt{max(0, ...)} en vez de arrastrarla al
camino final.
